The effect of weak bond disorder is especially delicate in two dimensions.
Harris's relevance criterion places the pure Ising model at the marginal
case, where the specific-heat exponent is zero and the criterion alone
does not decide the effect of randomness~\cite{Harris}.
Dotsenko and Dotsenko's fermionic analysis predicted a double logarithmic
specific-heat singularity~\cite{DD,DDReview}. Renormalization-group work
of Jug, Shalaev, Shankar, and Ludwig developed and refined the description
of logarithmic corrections to pure-model critical behavior
\cite{Jug,Shalaev,Shankar,Ludwig}. The predictions for different correlation
observables were not identical throughout this development. The relevant
physical picture here is marginal irrelevance: weak disorder decreases
under changes of scale, but only logarithmically. These works concern
thermodynamic quantities and correlations; control of a quenched law for
an entire interface also requires geometric estimates.

For the homogeneous critical model, the conformal description is
rigorous. Schramm's stochastic Loewner evolution supplies the candidate
interface laws~\cite{Schramm}. Smirnov established conformal convergence
of the FK-Ising fermionic observable~\cite{Smirnov}, and Chelkak and
Smirnov proved fermionic-observable universality for critical isoradial
Ising models~\cite{ChelkakSmirnov}. Chelkak, Duminil-Copin, Hongler,
Kemppainen, and Smirnov proved that the critical square-lattice spin
interface with Dobrushin boundary conditions converges to chordal
$\mathrm{SLE}_3$ in uniform distance modulo increasing
reparametrization~\cite[Theorem~1 and Equation~(1)]{CDHKS}.
This is the pure interface theorem underlying our comparison argument.

Rigorous universality results for deterministic interactions provide a
second point of comparison. Giuliani, Greenblatt, and Mastropietro proved
universality of infinite-plane energy-correlation scaling limits for
weak finite-range pair perturbations with lattice symmetries
\cite[Theorems~1.1--1.2]{GGM}. For non-integrable spin interfaces,
Greenblatt and Peltola constructed a Grassmann representation of a
discrete martingale and formulated the correlation estimates needed for
a proposed route to $\mathrm{SLE}_3$
\cite[Theorem~4.4 and Conjecture~4.5]{GP}. Their paper does not prove
interface convergence. For random bonds, Mahfouf proved quenched
FK-interface convergence to $\mathrm{SLE}_{16/3}$ in a regime where the
disorder strength tends to zero with the mesh
\cite[Theorem~1.6]{Mahfouf}. Av\'erous and Mahfouf established
high-probability crossing estimates for near-critical random-bond
FK models, also with vanishing disorder for $1<q\le4$
\cite[Theorems~1.5--1.6]{AM}. In the present theorem the nonzero disorder
strength is fixed as the mesh tends to zero, and the interface separates
spin clusters.

The main theorem of~\cite{R} treats symmetric two-valued bond disorder,
whereas its local scale transformation, moment estimates, and marginal
variance estimate are formulated for more general random local potentials.
We verify their hypotheses for the law~$\rho$ and its exact dual, and use
the pure and deterministic boundary-comparison results of~\cite{S}
through their formulations in~\cite{R}. The critical-point argument
adapts the torus method of~\cite[Section~8]{R} to match two separately
tuned trajectories within the same physical temperature family.
This matching replaces distributional self-duality.

Duality has a substantial history in the study of disordered critical
points. Chayes and Shtengel established vanishing magnetization and
critical correlation properties for self-dual disordered ferromagnets
under their stated hypotheses~\cite[Theorems~$1'$--$2'$]{ChayesShtengel}.
Our general one-edge law need not be self-dual. We use Kramers--Wannier
duality~\cite{KW} in the finite-torus form of Bugrij and Shadura, which
relates the four periodic/antiperiodic partition functions for arbitrary
positive inhomogeneous couplings~\cite[Equations~(4)--(5)]{BS}.
The torus comparison identifies the two tuned temperatures; a separate
argument identifies their common value with the spontaneous-magnetization
threshold.

The remaining argument combines classical probabilistic tools with
the geometric inputs of~\cite{R,S}. The Fortuin--Kasteleyn representation~\cite{FK}, the
Edwards--Sokal coupling~\cite{ES}, and the FKG lattice
inequality~\cite{FKG} relate spin barriers, random-cluster connections,
and boundary monotonicity. The decision-tree inequality of O'Donnell,
Saks, Schramm, and Servedio~\cite{OSSS}, extended to monotonic measures
by Duminil-Copin, Raoufi, and Tassion~\cite[Theorem~1.1]{DCRT}, is
applied at each fixed environment before averaging over the bonds.
The geometric use of crossing estimates to control random curves has
precedents in Aizenman--Burchard~\cite{AB} and
Kemppainen--Smirnov~\cite{KS}; Sheffield--Sun~\cite{SS} distinguish
convergence of driving functions from convergence in the uniform curve
topology. The specific passage from boundary barriers and trace
convergence to oriented-curve convergence used here is the progress
criterion of~\cite[Section~9]{R}.
